\section{Core performance}
\label{section:tailoring:core}

Before we continue, it is important to double-check whether we have used the correct
compiler settings. It makes no sense to benchmark a code which is not producing
tailored code for our particularly machinery (therefore it is important to
benchmark on the right system), and it makes no sense to benchmark a code that
is not using the latest compiler features.


First of all, we work with the the most aggressive optimisation that our
compiler can offer. The minimum choice is to compile with \texttt{-O3},
although some other compilers might offer further granularity. Also consider
using \texttt{-ffast-math} which enables aggressive floating-point
optimisation. Compiler documentation can be extensive, but there are often dedicated
pages which list optimisation flags.


Second, we create code that is tailored towards our particular machinery.
The correct setting for this depends once again on your compiler. 
With LLVM and GNU, you can use flags starting with \texttt{-mXXXX} to specify 
machinery-specific optimisations, i.e.~you can exactly say for which processor
generation, instruction set, \ldots~you want for the compiler to create binary code. Often we can tailor to the hardware with flags such as \texttt{-march=native}, with will automatically pickup the local architecture.
If your compile node is of the same type as your test node, then simply
use \texttt{-xhost} and the compiler will pick the highest instruction set
available.
The flag is called \texttt{-mhost} with GNU.

\begin{instructions}{\intelcompiler}
 With the Intel compiler, you sometimes can pick between the generic
 LLVM (community) versions of a feature and one optimised by Intel itself (with
 the latter obviously performing better on bespoke Intel architecture). For
 example, \texttt{-fopenmp} enables the general OpenMP implementation whereas
 \texttt{-fiopenmp} gives you an in-house version of the same. 
 
 And while \texttt{-xhost} uses Intel-specific optimisations, \texttt{-mhost}
 will use LLVM's optimisation stack, i.e.~Intel has replaced some optimisation
 steps with bespoke variants, which are enabled if and only if you go down the
 \texttt{-xhost} route. Be aware that Intel's optimisation passes typically will
 be disabled if you use the Intel toolchain on other vendor's hardware,
 i.e.~you'll get unoptimised code.
\end{instructions}

\begin{checklist}
 \item Your benchmark is compiled with the highest optimisation level (\texttt{-O3} plus a hardware-specific instruction set such as picked by \texttt{-xhost}).
\end{checklist}

\begin{instructions}{\maqao}
 If you run \maqao~over your code, the tool will report back to you what
 optimisations you might have missed out throughout the compilation.
  The top-level rubric \emph{Global} provides you with information on the optimisation chosen, but it also makes recommendations which options might lead to better runtime.

 The nice thing about \maqao~is that the tool can handle situations where you translate different translation units with different options.
 Also, the recommendation aspect goes beyond pure reporting and is helpful.
\end{instructions}

Not all compiler optimisations lead to decreases in runtime. With
tools like MAQAO it can be beneficial to \emph{experiment} with compiler
optimisations. For example, by default the GNU compilers do not implement loop
unrolling, and so MAQAO may initially suggest to switch on the
\texttt{-funroll-loops} flag. However, it is not a given that this will speed
up the code, and so we can the benchmark through MAQAO, with and without this
flag, to see if there is noticable differences in the total runtime, and the
runtimes of individual loops.

Some optimisations can implement aggressive optimisations. Many algorithms may
be robust to, for example, floating point approximations but it is worth
checking. If the code comes with test cases, or has known results which can be
checked against then it is very much advised to run these as validation checks.
Test the code with and without these aggressive optimisations to ensure that
these compiler flags have not noticably changed the code outcome. 

We recommend, if the analyst settles on using further compiler optimisations,
then these should be clearly documented in the performance report.  This is to
ensure that the code owner is aware of any aggressive optimisations that the
analyst uses, and the code owner makes any final decisions, or undertakes any
final checks, to ensure that the optimisations have not impacted the code
results.

\paragraph*{How it works}

\begin{itemize}[noitemsep]
  \item[$\boxplus$] Use the most aggressive compiler optimisation that preserve the code semantics (aka still return valid results) and nevertheless are tailored towards your assessment machinery.
  \item[$\boxplus$] Check what optimisations the users' build system uses by default. Often, this is a low-hanging fruit.
\end{itemize}


\paragraph*{How it does not work}

\begin{itemize}[noitemsep]
  \item[$\boxminus$] Use a low optimisation level and machine-generic optimisations and then draw conclusions on code behaviour on a particular machine.
  \item[$\boxminus$] Use aggresive mathematical approximations without validating correctness of the results.
\end{itemize}
